%  Origen: 03_esquema_solucion_alcance/entrega_2/tablas/esquema_solucion_alcance.xlsx
%  Hoja "3.2.2 Alcance de infraestructura del proyecto" (emplazamiento).

\begin{tablalafrox}{Modelo de emplazamiento híbrido: qué va dónde y por qué}%
  {tab:emplazamiento}%
  {L{4.0cm} L{3.3cm} Y}%
  {\cab{Componente} & \cab{Emplazamiento} & \cab{Justificación}}
  Recepción, picking y carga en los centros de distribución de Talca y
    Concepción
    & On-premise, con réplica a nube
    & Debe seguir operando durante un corte de enlace (restricción
      N.\textdegree\,2); el ambiente incluye montacargas y cámara de congelado
      a $-$22\,°C sin cobertura de señal. \\
  Plataformas de cross-docking de Curicó, Chillán y Los Ángeles
    & Gabinete de borde, con conectividad exclusiva por red móvil
    & Sin almacenamiento propio, ventana de operación de tres horas y sin fibra
      óptica disponible. \\
  Preventa y reparto en dispositivos de terreno
    & Aplicación que opera sin red, con sincronización diferida a nube
    & Deben operar un turno completo sin señal, en rutas rurales de hasta 240
      kilómetros (restricción N.\textdegree\,3). \\
  Trazabilidad de lote, cadena de frío, analítica, portal y componentes de
    inteligencia artificial
    & Nube pública, región Chile o Sudamérica
    & Volumen variable y sin acoplamiento a equipamiento físico; requiere
      elasticidad durante el peak de septiembre. \\
  Telemetría de los 42 camiones propios
    & Se integra desde la plataforma del proveedor externo que ya la provee
    & Fuente existente que se preserva; los 54 camiones de transportistas no
      están cubiertos y se incorporan mediante un mecanismo a definir con cada
      empresa (Decisión N.\textdegree\,5). \\
\end{tablalafrox}
